% ----- CHAUPAI 6 -----

\newpage

\subsection*{\deva{षष्ठ चौपाई} (Sixth Chaupai)}
\addcontentsline{toc}{subsection}{\deva{षष्ठ चौपाई} (Sixth Chaupai) — \deva{संकर सुवन...}}

\begin{minipage}[t]{0.55\textwidth}
\vspace{0pt}

{\large \linespread{1.5}\selectfont
\deva{संकर सुवन केसरी नंदन।}\\
\deva{तेज प्रताप महा जग बंदन॥}
\par}

\vspace{0.5em}

{\large \linespread{1.5}\selectfont
\telu{సంకర సువన కేసరీ నందన।}\\
\telu{తేజ ప్రతాప మహా జగ బందన॥}
\par}

\vspace{0.5em}

{\large \linespread{1.3}\selectfont
\textit{saṅkara suvana kesarī nandana |}\\
\textit{teja pratāpa mahā jaga bandana ||}
\par}
\end{minipage}%
\hfill
\begin{minipage}[t]{0.42\textwidth}
\vspace{0pt}
\begin{tcolorbox}[anchorbox]
\textbf{The Radiance of Grounded Service:} \textit{Shankara} represents the highest spiritual intellect and ego-less discipline, while \textit{Kesari} represents a highly practical, grounded, forest-dwelling leader. Hanuman embodies the perfect balance of both: he possesses the brilliant mind of a supreme thinker, yet uses his fierce protective energy (\textit{Teja Pratapa}) entirely for the selfless service of others rather than personal gain. The world reveres him (\textit{Jaga Bandana}) because he proves that true majestic radiance comes from balancing high-minded vision with practical, ego-less dedication to your community.
\vspace{0.5em}\newline Jambavan honors his earthly, grounded origins, proving true radiance stems from one's actions, not just their birth.\newline\null\hfill\href{https://www.valmiki.iitk.ac.in/sloka?field_kanda_tid=4&language=dv&field_sarga_value=66}{\textit{Kishkindha Kanda, 66:29}}
\end{tcolorbox}
\end{minipage}

\vspace{0.5em}
\subsubsection*{\textbf{Visual Rhythm Map:}}
\begin{table}[h!]
\centering
\renewcommand{\arraystretch}{1.2}
\setlength{\tabcolsep}{0pt}
\begin{tabularx}{\textwidth}{|*{16}{>{\centering\arraybackslash\hsize=1.0\hsize}X|}}
\hline
\multicolumn{4}{|>{\centering\arraybackslash\hsize=4.0\hsize}X|}{\textbf{\guru{सं}\laghu{क}\laghu{र}} \par (4)} &
\multicolumn{3}{>{\centering\arraybackslash\hsize=3.0\hsize}X|}{\textbf{\laghu{सु}\laghu{व}\laghu{न}} \par (3)} &
\multicolumn{5}{>{\centering\arraybackslash\hsize=5.0\hsize}X|}{\textbf{\guru{के}\laghu{स}\guru{री}} \par (5)} &
\multicolumn{4}{>{\centering\arraybackslash\hsize=4.0\hsize}X|}{\textbf{\guru{नं}\laghu{द}\laghu{न}} \par (4)} \\
\hline
\end{tabularx}
\vspace{-1.5em}
\end{table}

\begin{table}[h!]
\centering
\renewcommand{\arraystretch}{1.2}
\setlength{\tabcolsep}{0pt}
\begin{tabularx}{\textwidth}{|*{16}{>{\centering\arraybackslash\hsize=1.0\hsize}X|}}
\hline
\multicolumn{3}{|>{\centering\arraybackslash\hsize=3.0\hsize}X|}{\textbf{\guru{ते}\laghu{ज}} \par (3)} &
\multicolumn{4}{>{\centering\arraybackslash\hsize=4.0\hsize}X|}{\textbf{\laghu{प्र}\guru{ता}\laghu{प}} \par (4)} &
\multicolumn{3}{>{\centering\arraybackslash\hsize=3.0\hsize}X|}{\textbf{\laghu{म}\guru{हा}} \par (3)} &
\multicolumn{2}{>{\centering\arraybackslash\hsize=2.0\hsize}X|}{\textbf{\laghu{ज}\laghu{ग}} \par (2)} &
\multicolumn{4}{>{\centering\arraybackslash\hsize=4.0\hsize}X|}{\textbf{\guru{बं}\laghu{द}\laghu{न}} \par (4)} \\
\hline
\end{tabularx}
\vspace{-1.5em}
\end{table}

\vspace{1em}
\textbf{ \deva{सरल-संस्कृतम्}:}\\
 \deva{भवान् शङ्करस्य अवतारः, केसरी-पुत्रः च असि।}\\
 \deva{भवतः तेजः प्रतापः च महान् अस्ति, सम्पूर्णं जगत् भवन्तं वन्दते (नमति)॥}\\


You are the incarnation of Shiva and the son of Kesari.\\
Your radiance and valor are immense, and you are revered by the whole universe.


\vspace{1em}
\newpage
\textbf{\deva{प्रतिपदार्थः} (Meaning of each word):}
\begin{table}[h!]
\centering
\renewcommand{\arraystretch}{1.2}
\begin{tabularx}{\textwidth}{@{} l l X X @{}}
\toprule
\textbf{Awadhi} & \textbf{Sanskrit} & \textbf{English} & \textbf{Grammar \& Notes} \\
\midrule
\deva{संकर} \gotchaicon / \telu{సంకర} / \textit{saṅkara}  &  \deva{शङ्करस्य}  &  Of Shiva  & \\ 
\deva{संकर सुवन} / \telu{సంకర సువన} / \textit{saṅkara suvana}  &  \deva{शङ्करस्य पुत्रः (हनुमान्)}  &  Son of Shiva  &   \\
\deva{केसरी नंदन} / \telu{కేసరీ నందన} / \textit{kesarī nandana}  &  \deva{केसरिणः नन्दनः}  &  Son of Kesari  & \\ 
\deva{तेज} / \telu{తేజ} / \textit{teja}  &  \deva{तेजस्वित्वम्}  &  Radiance / Energy  &   \\
\deva{प्रताप} / \telu{ప్రతాప} / \textit{pratāpa}  &  \deva{प्रतापः}  &  Valor / Prowess  &   \\
\deva{महा} / \telu{మహా} / \textit{mahā}  &  \deva{महान्}  &  Great  &   \\
\deva{जग} / \telu{జగ} / \textit{jaga} & \deva{जगति} & In the world & \\
\deva{बंदन} / \telu{బందన} / \textit{bandana} & \deva{वन्दनीयः} & Adored / Worshiped & \\
 
\bottomrule
\end{tabularx}
\end{table}

\begin{tcolorbox}[gotchabox]
\begin{itemize}
    \item \textbf{Phonetic Shifts:} Sanskrit \deva{शङ्कर} (Shankara) seamlessly shifts into the Awadhi \deva{संकर} (Sankara) losing the palatal 'sh'. Likewise, \deva{वन्दन} (Vandana) hardens into \deva{बंदन} (Bandana) with the \deva{व} $\rightarrow$ \deva{ब} shift.
    \item \textbf{The Secret of \deva{सुवन} (The Glide / \deva{श्रुति}):} If \deva{वन्दन} becomes \deva{बंदन}, why didn't \deva{सुवन} become \deva{सुबन}? The word \deva{सुवन} comes from the Sanskrit root \deva{सू} (to beget) plus the noun suffix \deva{अन}. Saying \deva{सू} and \deva{अन} rapidly together naturally creates a soft 'w/v' sound as a breath-bridge (a \textit{Glide} or \deva{व-श्रुति}). Because this \deva{व} is just a fluid phonetic bridge between two vowels—and not a heavy structural consonant—the Awadhi speakers kept it soft instead of hammering it into a hard \deva{ब}.
\end{itemize}
\end{tcolorbox}

\newpage
